\documentclass[11pt,a4paper]{article}
	
	\usepackage{amsmath,amssymb,amsthm,mathtools,amsfonts}
	\usepackage{breqn}
	\usepackage[utf8]{inputenc}
	\usepackage[danish]{babel}
	\usepackage[T1]{fontenc}
	
	%%%%%%%%%%%%%%%%%%%%%%%%%%%%%%%%%%%%%%%%%%%%%%%%%%%%%%%%%%%%%%%%%%%%%%%%%%%%%% Fonts
	
	\usepackage{fontspec}
	\setmainfont[Numbers=OldStyle]{EB Garamond}
	\newfontfamily\plantinc[Numbers=OldStyle]{Plantin MT Pro Bold Condensed}
	\newfontfamily\plantinsb{Plantin MT Pro Semibold}
	\newfontfamily\garamond[Numbers=OldStyle]{EB Garamond}
	\newfontfamily\plantin[Numbers=OldStyle]{Plantin MT Pro}
	\newfontfamily\wal{Walbaum Com}
	\newfontfamily\klav{Klavika}
	\newfontfamily\klavl{Klavika light}
	\newfontfamily\quotefont[Ligatures=TeX]{Linux Libertine O}
	
	%%%%%%%%%%%%%%%%%%%%%%%%%%%%%%%%%%%%%%%%%%%%%%%%%%%%%%%%%%%%%%%%%%%%%%%%%%%%%% Colors
	
	\usepackage[svgnames]{xcolor}
	
	\definecolor{hgblue}{RGB}{10, 40, 80}
	\definecolor{hgred}{RGB}{140, 10, 10}
	\definecolor{hggrey}{RGB}{215, 215, 210}
	\definecolor{hggreen}{RGB}{145, 205, 205}
	
	%%%%%%%%%%%%%%%%%%%%%%%%%%%%%%%%%%%%%%%%%%%%%%%%%%%%%%%%%%%%%%%%%%%%%%%%%%%%%% Other Graphics
	
	\usepackage{graphicx}
	\graphicspath{{./img/}}
	\usepackage{tikz}
	\usepackage{placeins}
	\usepackage[modulo]{lineno}
	%\usepackage{geometry}
	
	\usepackage[pdfborder={0 0 0}]{hyperref} %% Ingen firkant rundt om referencer

	\usepackage{multicol}
	\usepackage{cellspace}
		\setlength\cellspacetoplimit{100pt}
		\setlength\cellspacebottomlimit{100pt}
	\usepackage{tabularx}
		\newcolumntype{b}{>{\hsize=\hsize}>{\centering\arraybackslash}X}
	\usepackage{siunitx}
	\usepackage{cellspace}
		\setlength\cellspacetoplimit{4pt}
		\setlength\cellspacebottomlimit{4pt}
	
	\usepackage{enumitem}% better controls with enumerating
	\usepackage{wrapfig}
		%\setlength{\wrapoverhang}{\marginparwidth}
		%\addtolength{\wrapoverhang}{\marginparsep}
		\setlength{\columnsep}{0pt}
		\setlength{\intextsep}{-10pt}
	\usepackage{caption}
	%\usepackage{dirtytalk}
	%\usepackage{tasks}
	
	

	%\setlength{\parskip}{0.5\baselineskip}
	%\setlength{\parindent}{0cm}
	
	\setlist[enumerate]{label=\alph*),left=20pt}
	%\setlist[enumerate]{leftmargin=\parindent}
	%\setlist[enumerate]{nosep}
	
	\usepackage{lipsum}
	\usepackage{comment}
	\usepackage{ulem}
	
	%\reversemarginpar
	
	%%%%%%%%%%%%%%%%%%%%%%%%%%%%%%%%%%%%%%%%%%%%%%%%%%%%%%%%%%%%%%%%%%%%%%%%%%%%%% New commands
	
	\newcommand\svar[1]{\newline\textcolor{red}{\textbf{Svar}\\#1}}
	
	\usepackage[h]{esvect}
	\newcommand{\twvect}[2]{\ensuremath{\begin{pmatrix}#1\\#2\end{pmatrix}}}
	
	\newcommand*\quotesize{20}
	\newcommand*{\openquote}
		{\tikz[remember picture,overlay,xshift=-3ex,yshift=-0ex]
			\node (OQ) {\quotefont\fontsize{\quotesize}{\quotesize}\selectfont``};\kern0pt}

	\newcommand*{\closequote}[1]
		{\tikz[remember picture,overlay,xshift=2ex,yshift={#1}]
			\node (CQ) {\quotefont\fontsize{\quotesize}{\quotesize}\selectfont''};}
			
	\newcommand{\qm}[1]{‘‘#1”}
	\newcommand{\iquote}[1]{\begin{quote}\itshape \openquote#1\hfill\closequote{0ex} \end{quote}}

	\newcommand{\source}[1]{\footnote{Source: \href{#1}{#1}}}
	
	\newcommand{\glos}[3]{\dotuline{#1}\marginpar{\scriptsize\textit{#3} #2}}
	
	%%%%%%%%%%%%%%%%%%%%%%%%%%%%%%%%%%%%%%%%%%%%%%%%%%%%%%%%%%%%%%%%%%%%%%%%%%%%%% Sections
	
	\usepackage{titlesec}
	\titleformat{\subsection}{\color{hgred}\plantinc\large}{}{}{}
	
	\usepackage{titling}
	\setlength{\droptitle}{-12ex}
	\pretitle{\begin{flushleft}\plantinc\huge\color{hgblue}}
	\posttitle{\par\end{flushleft}}
	\preauthor{\begin{flushleft}\Large}
	\postauthor{\end{flushleft}}
	\predate{\begin{flushleft}}
	\postdate{\end{flushleft}}
	
	\setcounter{secnumdepth}{0}
	
	%%%%%%%%%%%%%%%%%%%%%%%%%%%%%%%%%%%%%%%%%%%%%%%%%%%%%%%%%%%%%%%%%%%%%%%%%%%%%% Header & Footer
	
	\usepackage{bigfoot}
	\DeclareNewFootnote{a}
		\renewcommand{\thefootnotea}{\fnsymbol{footnotea}}
			\newcommand{\footnotes}[1]{\footnotea[2]{#1}} % here you can specify what symbol you want in footnotes
			\newcommand{\footnoted}[1]{\footnotea[3]{#1}} % for a second footnote on a page
			\MakePerPage{footnotes}
	
			%%%%%%%%%%%%%%%%%%%%%%%%%%%%%%%%%%%%%%%%%%%%%%%%%%%%%%%%%%%%%%%%%%%%%%%%%%%%%% Document

\begin{document}
%\pagecolor{FloralWhite}
\title{The Black Cat}
%\date{\vspace{-3em}} % I tilfældet ingen dato
\date{1843}
\author{Edgar Allen Poe}
\subsection{Introduction}
The Black Cat is a story of an unnamed man who starts out by describing how he was very caring as a child and loved animals. Later in life, he and his wife acquire a lot of pets and among these a black cat named Pluto, which he becomes very fond of. As the man grows older, his character changes for the worse, and he has trouble controlling his temper (it is indicated that this change may be caused by alcohol abuse). One day he becomes frustrated with the cat as it appears to have become frightened of him, and he loses his temper and hurts the cat so badly that it loses an eye. With only one eye, the cat frustrates him even more, and he ends up hanging it in the basement. The narrator’s house mysteriously burns down that same night. Soon after that, the man starts to feel remorse for having killed the cat and looks to replace it. One day, he comes across a very similar black cat and decides to take it home with him. The next morning he discovers that his new cat also has only one eye. 

The excerpt that follows is the second half of the story.
\maketitle
    
\thispagestyle{plain}
\setcounter{page}{1}
\linenumbers
\noindent […]\\[5pt]

\noindent It followed my footsteps with a pertinacity which it would be difficult to make the reader comprehend. Whenever I sat, it would crouch beneath my chair, or spring upon my knees, covering me with its loathsome caresses. If I arose to walk, it would get between my feet and thus nearly throw me down, or, fastening its long and sharp claws in my dress, clamber, in this manner, to my breast. At such times, although I longed to destroy it with a blow, I was yet withheld from so doing, partly by a memory of my former crime, but chiefly -- let me confess it at once -- by absolute dread of the beast.

This dread was not exactly a dread of physical evil -- and yet I should be at a loss how otherwise to define it. I am almost ashamed to own -- yes, even in this felon's cell, I am almost ashamed to own -- that the terror and horror with which the animal inspired me, had been heightened by one of the merest chimeras it would be possible to conceive. My wife had called my attention, more than once, to the character of the mark of white hair, of which I have spoken, and which constituted the sole visible difference between the strange beast and the one I had destroyed. The reader will remember that this mark, although large, had been originally very indefinite; but, by slow degrees -- degrees nearly imperceptible, and which for a long time my Reason struggled to reject as fanciful -- it had, at length, assumed a rigorous distinctness of outline. It was now the representation of an object that I shudder to name -- and for this, above all, I loathed, and dreaded, and would have rid myself of the monster had I dared -- it was now, I say, the image of a hideous -- of a ghastly thing -- of the GALLOWS ! […]

Neither by day nor by night knew I the blessing of Rest any more! During the former the creature left me no moment alone; and, in the latter, I started, hourly, from dreams of unutterable fear, to find the hot breath of the thing upon my face, and its vast weight -- an incarnate Night-Mare that I had no power to shake off -- incumbent eternally upon my heart!

Beneath the pressure of torments such as these, the feeble remnant of the good within me succumbed. Evil thoughts became my sole intimates -- the darkest and most evil of thoughts. The moodiness of my usual temper increased to hatred of all things and of all mankind; while, from the sudden, frequent, and ungovernable outbursts of a fury to which I now blindly abandoned myself, my uncomplaining wife, alas! was the most usual and the most patient of sufferers.

One day she accompanied me, upon some household errand, into the cellar of the old building which our poverty compelled us to inhabit. The cat followed me down the steep stairs, and, nearly throwing me headlong, exasperated me to madness. Uplifting an axe, and forgetting, in my wrath, the childish dread which had hitherto stayed my hand, I aimed a blow at the animal which, of course, would have proved instantly fatal had it descended as I wished. But this blow was arrested by the hand of my wife. Goaded, by the interference, into a rage more than demoniacal, I withdrew my arm from her grasp and buried the axe in her brain. She fell dead upon the spot, without a groan.

This hideous murder accomplished, I set myself forthwith, and with entire deliberation, to the task of concealing the body. I knew that I could not remove it from the house, either by day or by night, without the risk of being observed by the neighbors. Many projects entered my mind. […] I determined to wall it up in the cellar -- as the monks of the middle ages are recorded to have walled up their victims.

For a purpose such as this the cellar was well adapted. Its walls were loosely constructed, and […] in one of the walls was a projection, caused by a false chimney, or fireplace, that had been filled up, and made to resemble the rest of the cellar. I made no doubt that I could readily displace the bricks at this point, insert the corpse, and wall the whole up as before, so that no eye could detect anything suspicious.

And in this calculation I was not deceived. By means of a crowbar I easily dislodged the bricks, and, having carefully deposited the body against the inner wall, I propped it in that position, while, with little trouble, I re-laid the whole structure as it originally stood. […] When I had finished, I felt satisfied that all was right. The wall did not present the slightest appearance of having been disturbed. The rubbish on the floor was picked up with the minutest care. I looked around triumphantly, and said to myself -- "Here at least, then, my labor has not been in vain."

My next step was to look for the beast which had been the cause of so much wretchedness; for I had, at length, firmly resolved to put it to death. Had I been able to meet with it, at the moment, there could have been no doubt of its fate; but it appeared that the crafty animal had been alarmed at the violence of my previous anger, and forebore to present itself in my present mood. It is impossible to describe, or to imagine, the deep, the blissful sense of relief which the absence of the detested creature occasioned in my bosom. It did not make its appearance during the night -- and thus for one night at least, since its introduction into the house, I soundly and tranquilly slept; aye, slept even with the burden of murder upon my soul!

The second and the third day passed, and still my tormentor came not. Once again I breathed as a freeman. The monster, in terror, had fled the premises forever! I should behold it no more! My happiness was supreme! The guilt of my dark deed disturbed me but little. Some few inquiries had been made, but these had been readily answered. Even a search had been instituted -- but of course nothing was to be discovered. I looked upon my future felicity as secured.

Upon the fourth day of the assassination, a party of the police came, very unexpectedly, into the house, and proceeded again to make rigorous investigation of the premises. Secure, however, in the inscrutability of my place of concealment, I felt no embarrassment whatever. […] The police were thoroughly satisfied and prepared to depart. The glee at my heart was too strong to be restrained. I burned to say if but one word, by way of triumph, and to render doubly sure their assurance of my guiltlessness.

"Gentlemen," I said at last, as the party ascended the steps, "I delight to have allayed your suspicions. I wish you all health, and a little more courtesy. By the bye, gentlemen, this -- this is a very well constructed house." […] and here, through the mere phrenzy of bravado, I rapped heavily, with a cane which I held in my hand, upon that very portion of the brick-work behind which stood the corpse of the wife of my bosom.

[…] I was answered by a voice from within the tomb! -- by a cry, at first muffled and broken, like the sobbing of a child, and then quickly swelling into one long, loud, and continuous scream, utterly anomalous and inhuman -- a howl -- a wailing shriek, half of horror and half of triumph, such as might have arisen only out of hell […] 

Of my own thoughts it is folly to speak. Swooning, I staggered to the opposite wall. For one instant the party upon the stairs remained motionless, through extremity of terror and of awe. In the next, a dozen stout arms were toiling at the wall. It fell bodily. The corpse, already greatly decayed and clotted with gore, stood erect before the eyes of the spectators. Upon its head, with red extended mouth and solitary eye of fire, sat the hideous beast whose craft had seduced me into murder, and whose informing voice had consigned me to the hangman. I had walled the monster up within the tomb!








































\end{document}